%% ============================================================
%%  preamble.tex – Pakete & Konfiguration
%%  Hochschule Heilbronn | Bachelor-Thesis
%% ============================================================

%% ---- Sprache & Encoding ------------------------------------
\usepackage[utf8]{inputenc}
\usepackage[T1]{fontenc}
\usepackage[ngerman]{babel}
\usepackage{csquotes}      
\usepackage{physics}

% \enquote{} für korrekte Anführungszeichen

%% ---- Schrift -----------------------------------------------
\usepackage{lmodern}
\usepackage{microtype}            % Besserer Textsatz, weniger Overfull hboxes

%% ---- Seitenlayout ------------------------------------------
%  HHN-Norm: links 3cm (Bindung), rechts 2.5cm, oben/unten 2.5cm
\usepackage[
  a4paper,
  left=3cm,
  right=2.5cm,
  top=2.5cm,
  bottom=2.5cm,
  bindingoffset=5mm
]{geometry}

%% ---- Zeilenabstand -----------------------------------------
\usepackage{setspace}
\onehalfspacing

%% ---- Kopf- & Fußzeilen (KOMA) ------------------------------
\usepackage[
  headsepline,
  automark
]{scrlayer-scrpage}
\pagestyle{scrheadings}
\automark[section]{chapter}
\clearpairofpagestyles
\ihead{\headmark}
\ohead{\pagemark}
\cfoot{}
% Kapitelseiten erhalten dieselbe Kopfzeile
\renewcommand*{\chapterpagestyle}{scrheadings}

%% ---- Grafik & Abbildungen ----------------------------------
\usepackage{graphicx}
\graphicspath{{figures/}}
\usepackage{float}
\usepackage{caption}
\captionsetup{
  font=small,
  labelfont=bf,
  format=hang
}
\usepackage{subcaption}
\usepackage{multirow}


%% ---- Tabellen ----------------------------------------------
\usepackage{booktabs}             % \toprule, \midrule, \bottomrule
\usepackage{tabularx}
\usepackage{longtable}
\usepackage{array}

%% ---- Mathematik --------------------------------------------
\usepackage{amsmath}
\usepackage{amssymb}
\usepackage{eurosym}
%% ---- Quellcode (Listings) ----------------------------------
\usepackage{listings}
\usepackage{xcolor}
\usepackage{graphicx}
\usepackage{minted} 

\definecolor{lstbg}{gray}{0.97}
\definecolor{lstcomment}{rgb}{0.13,0.55,0.13}
\definecolor{lstkeyword}{rgb}{0.0,0.0,0.75}
\definecolor{lststring}{rgb}{0.63,0.13,0.13}

\lstset{
  backgroundcolor=\color{lstbg},
  basicstyle=\ttfamily\small,
  keywordstyle=\color{lstkeyword}\bfseries,
  commentstyle=\color{lstcomment}\itshape,
  stringstyle=\color{lststring},
  numbers=left,
  numberstyle=\tiny\color{gray},
  numbersep=8pt,
  frame=single,
  framerule=0.4pt,
  rulecolor=\color{gray!40},
  breaklines=true,
  breakatwhitespace=true,
  captionpos=b,
  tabsize=2,
  showstringspaces=false
}

%% ---- Abkürzungsverzeichnis ---------------------------------
\usepackage{acronym}

%% ---- Literaturverwaltung (BibLaTeX + Biber) ----------------
\usepackage[
  backend=biber,
  style=ieee,       % Autor-Jahr-Stil (Harvard)
  % style=numeric         % Nummern-Stil [1], [2], ...
  % style=ieee            % IEEE-Stil
  sorting=nyt,
  maxbibnames=99,
  maxcitenames=2,
  giveninits=true,
  uniquename=init
]{biblatex}
\addbibresource{literatur.bib}

%% ---- Hyperlinks & PDF-Metadaten ----------------------------
\usepackage[
  colorlinks=true,
  linkcolor=black,
  citecolor=black,
  urlcolor=blue,
  bookmarksnumbered=true,
  pdfborder={0 0 0}
]{hyperref}
\usepackage{bookmark}

%% ---- Hilfspackages -----------------------------------------
\usepackage{lipsum}               % Blindtext zum Testen → später entfernen


%% --- Code-----------------------------------------------------


\lstdefinestyle{python}{
  language=Python,
  basicstyle=\ttfamily\small,
  keywordstyle=\color{blue!70!black},
  stringstyle=\color{green!40!black},
  commentstyle=\color{gray}\itshape,
  numbers=left,
  numberstyle=\tiny\color{gray},
  numbersep=8pt,
  frame=single,
  rulecolor=\color{gray!40},
  backgroundcolor=\color{gray!5},
  breaklines=true,
  showstringspaces=false,
  captionpos=b,
  extendedchars=true,
  literate={ä}{{\"a}}1 {ö}{{\"o}}1 {ü}{{\"u}}1 {Ä}{{\"A}}1 {Ö}{{\"O}}1 {Ü}{{\"U}}1 {ß}{{\ss}}1
}
\lstset{style=python}